\section{Workflow Description}

The project followed a sequential workflow consisting of data ingestion,
inspection, cleaning, exploratory analysis, visualization, and interpretation.
This structure was chosen to make the analytical process transparent and
repeatable. Danchev describes reproducibility as a principle that should be
incorporated throughout the data-science lifecycle rather than treated only as
a final documentation activity \parencite{danchev2022}.

\subsection{Ingestion}

The CSV file was stored within the repository and loaded with
\texttt{pandas.read\_csv()}. The notebook then inspected the dimensions, column
names, data types, missing values, duplicated rows, categorical values, and
numerical ranges before applying any transformations.

\subsection{Cleaning}

Two reusable cleaning functions were created. The first converted
\texttt{last\_review} to a datetime type and handled missing review-rate values.
The second removed records that violated explicit validity rules, including
non-positive prices.

\subsection{Exploratory Analysis}

A reusable exploratory-analysis function generated numerical summaries,
neighbourhood-level summaries, room-type summaries, combined neighbourhood
and room-type comparisons, and a Pearson correlation matrix. Additional
analyses examined zero availability and compared mean and median prices before
and after limiting prices to the 99th percentile.

\subsection{Visualizations}

Four visualizations were created:

\begin{enumerate}[label=\arabic*.]
    \item median price by neighbourhood group and room type;
    \item price distributions by room type;
    \item median annual availability by neighbourhood group and room type;
    \item correlations between numerical variables.
\end{enumerate}

\subsection{Summary}

The final notebook section interpreted the calculated results, documented the
main limitations and assumptions, and identified questions that could not be
resolved from the available variables.